\documentclass[11pt]{article}
\usepackage[utf8]{inputenc}
\usepackage[T1]{fontenc}
\usepackage{lmodern}
\usepackage[margin=1in]{geometry}
\usepackage{amsmath,amssymb,amsthm,mathtools}
\usepackage{hyperref}
\usepackage[capitalise,noabbrev]{cleveref}

\theoremstyle{plain}
\newtheorem{theorem}{Theorem}[section]
\newtheorem{proposition}[theorem]{Proposition}
\newtheorem{lemma}[theorem]{Lemma}
\newtheorem{corollary}[theorem]{Corollary}
\newtheorem{conjecture}[theorem]{Conjecture}
\theoremstyle{definition}
\newtheorem{definition}[theorem]{Definition}
\newtheorem{example}[theorem]{Example}
\theoremstyle{remark}
\newtheorem{remark}[theorem]{Remark}

\title{{{TITLE}}}
\author{{{AUTHOR}}}
\date{\today}

\begin{document}
\maketitle

\begin{abstract}
Write the abstract here.
\end{abstract}

\section{Introduction}\label{sec:intro}
State the main result early.

\begin{theorem}\label{thm:main}
Let $X$ be \dots. Then \dots.
\end{theorem}

By \cref{thm:main} \dots, as in \cite{example2024}.

\section{Preliminaries}

\section{Proof of the main theorem}
\begin{proof}[Proof of \cref{thm:main}]
\dots
\end{proof}

\bibliographystyle{alpha}
\bibliography{references}
\end{document}
